% ECONOMICS_PROBLEM: {"id": "E32-471", "title": "Welfare costs of fluctuations", "jel_code": "E32", "source_id": "OP-0237"}
% !TeX program = xelatex
\documentclass[11pt,a4paper]{article}
\usepackage[margin=23mm]{geometry}
\usepackage{fontspec}
\setmainfont{Latin Modern Roman}
\usepackage{amsmath,amssymb,mathtools}
\usepackage{xcolor,enumitem,booktabs,longtable,array,xurl,hyperref}
\definecolor{muted}{HTML}{53636C}
\hypersetup{colorlinks=true,urlcolor=blue,linkcolor=blue}
\setlength{\parindent}{0pt}
\setlength{\parskip}{5pt}
\newcommand{\R}{\mathbb R}
\newcommand{\E}{\mathbb E}
\newcommand{\N}{\mathbb N}
\newcommand{\ind}{\mathbf 1}
\newcommand{\Var}{\operatorname{Var}}
\newcommand{\Cov}{\operatorname{Cov}}
\newcommand{\supp}{\operatorname{supp}}
\DeclareMathOperator*{\argmin}{arg\,min}
\DeclareMathOperator*{\argmax}{arg\,max}
\newcommand{\entrymeta}[3]{{\sffamily\small\color{muted}\textbf{\texttt{#1}}\quad|\quad #2\par #3\par}}
\newcommand{\Problemref}[1]{\texttt{#1}}
\newcommand{\JELref}[2]{\texttt{#2} (JEL \texttt{#1})}
\begin{document}
\section*{Welfare costs of fluctuations}
\textbf{Problem ID:} \texttt{E32-471}\quad \textbf{JEL:} \texttt{E32}\par
% BEGIN ECONOMICS STATEMENT
\label{problem:OP-0237}
{\sffamily\small\color{muted}Original subject group: Business cycles and macro dynamics\par}
\label{definitions:problem:30007518}
\textbf{Audit classification:} Background; proposed extension. Survey willingness-to-pay findings do not establish the exact utility-identification problem as open. Verification concerns the cited version, not first priority or proof that this exact editorial specification remains unsolved on October 8, 2026.
\subsection*{Definitions and precise research target}
\paragraph{Editorial mathematical specification.} Fix a population of household types \(i\), social weights \(\omega_i\ge0\) summing to one, discount factor \(\beta\in(0,1)\), and utility \(u_i(c,\ell)\) over consumption and labor. Let \((c_{it},\ell_{it})\) be allocations in the observed economy and \((\bar c_{it},\bar\ell_{it})\) those in a specified counterfactual that removes aggregate innovations while retaining specified idiosyncratic risks and policy rules and recomputing equilibrium resources. Assume utility is increasing in consumption, discounted expectations are finite, and any compensation exists uniquely with \(\lambda_i>-1\). Define each type’s consumption-equivalent fluctuation cost \(\lambda_i\) by
\[
E\sum_{t\ge0}\beta^t u_i((1+\lambda_i)c_{it},\ell_{it})
=E\sum_{t\ge0}\beta^t u_i(\bar c_{it},\bar\ell_{it}).
\]
Identify or bound the distribution of \(\lambda_i\) and its socially weighted mean from household consumption, employment, beliefs, and elicited willingness to pay. Explicitly distinguish utility-based compensation from hypothetical survey payments: specify the assumptions under which reported payments identify \(\lambda_i\), and test framing, perceived feasibility, and income effects. The stabilization counterfactual must incorporate changes in unemployment risk and general-equilibrium prices rather than simply smoothing measured consumption ex post. Assess sensitivity to risk aversion, persistent uncertainty, incomplete insurance, and welfare weights. The exercise yields costs conditional on declared preferences and stabilization technology, not a preference-free universal welfare cost of cycles.

An implementable stabilization economy must satisfy its own resource constraints; removing aggregate innovations does not generally preserve expected resources. State separately whether the counterfactual is a feasible equilibrium policy or a statistical consumption-smoothing benchmark. Let \(V_i(\lambda)=E\sum_t\beta^tu_i((1+\lambda)c_{it},\ell_{it})\); strict monotonicity of \(V_i\), continuity, and the counterfactual value lying in its range establish uniqueness, rather than merely assuming it. The weighted mean \(\sum_i\omega_i\lambda_i\) is a descriptive average of individual consumption equivalents and generally is not a representative social compensation. The latter solves \(\sum_i\omega_iV_i(\lambda^{\rm soc})=\sum_i\omega_i\bar V_i\). Survey payments identify either quantity only under an explicit elicitation and perception model.
\subsection*{Variants and assumption changes}
The following are \emph{editorial variants}, not additional published open conjectures.
\begin{enumerate}
\item \emph{CRRA consumption benchmark.} Fix \(u(c)=c^{1-\gamma}/(1-\gamma)\), \(\gamma>0\), \(\gamma\ne1\), and exogenous labor. Compare consumption-equivalent costs under transient versus permanent aggregate risk, holding means and idiosyncratic innovations fixed within this statistical benchmark.
\item \emph{Feasible insurance policy.} Replace perfect aggregate smoothing by a budget-feasible unemployment-insurance policy with endogenous labor and tax financing. Calculate household and social consumption equivalents relative to the actual policy, including the insurance and incentive effects.
\end{enumerate}
\subsection*{References and source audit}
\begin{itemize}
\item Dimitris Georgarakos; Kwang Hwan Kim; Olivier Coibion; Myungkyu Shim; Myunghwan Andrew Lee; Yuriy Gorodnichenko; Geoff Kenny; Seowoo Han; Michael Weber. \emph{How Costly Are Business Cycle Volatility and Inflation? A Vox Populi Approach}. 2025. Locator: IZA DP17675 metadata and abstract, February 2025; abstract supplies results. Primary record and abstract/summary checked; no fulltext claim is made. \url{https://www.iza.org/en/publications/dp/17675/how-costly-are-business-cycle-volatility-and-inflation-a-vox-populi-approach}.
\end{itemize}
Survey willingness-to-pay findings do not establish the exact utility-identification problem as open.
\begin{enumerate}\raggedright
\item\hypertarget{definitions:source:30007518:1}{} Dimitris Georgarakos; Kwang Hwan Kim; Olivier Coibion; Myungkyu Shim; Myunghwan Andrew Lee; Yuriy Gorodnichenko; Geoff Kenny; Seowoo Han; Michael Weber. 2025. \emph{How Costly Are Business Cycle Volatility and Inflation? A Vox Populi Approach}. IZA DP17675 metadata and abstract, February2025; abstract supplies results.
\url{https://www.iza.org/en/publications/dp/17675/how-costly-are-business-cycle-volatility-and-inflation-a-vox-populi-approach}.
DOI: \url{https://doi.org/10.3386/w33476}.
\end{enumerate}

\subsection*{Common conventions: Source mathematical conventions}

These conventions apply to each entry unless explicitly replaced.

\textbf{Models and identification.} A primitive model \(m\) specifies all probability laws, feasible choices, information, equilibrium and measurement rules used by its target. The admissible class \(\mathcal M\) is fixed independently of outcomes. If \(P_O(m)\) is its observable probability law and \(\tau(m)\) the target, its identified set at observable law \(P\) is
\[
 \mathcal I_\tau(P)=\{\tau(m):m\in\mathcal M,\ P_O(m)=P\}.
\]
Point identification means that this set is a singleton. An empty set signals incompatibility of data and assumptions. Coordinate bounds need not describe the joint set. Bounds are sharp only if every retained target is attainable or arbitrarily approachable by compatible models. Finite bounds require explicit restrictions on unobserved quantities; unrestricted confounding, hidden wealth, extrapolation or arbitrary equilibrium selection can yield uninformative sets.

\textbf{Causal interventions.} A potential outcome \(Y_i(p)\) is the outcome under a complete policy regime \(p\), including timing, financing, information and market spillovers. The target population and units are fixed before treatment. With interference, use the complete assignment vector or a justified exposure mapping; individual no-interference notation alone cannot identify an economy-wide intervention. A difference of mean potential outcomes does not require the cross-world joint distribution. Conditioning on post-treatment migration, survival, enrollment or employment changes the estimand and needs separate justification.

\textbf{Statistical statements.} An estimator, confidence set, forecast or test specifies a sampling law, model class and loss/coverage criterion. Uniform \((1-\alpha)\)-coverage means \(\inf_{m\in\mathcal M}\Pr_m\{\tau(m)\in C_n\}\geq1-\alpha\), or an explicitly asymptotic version. The whole parameter space trivially has coverage, so informativeness requires a stated width, risk or power objective. A forecasting improvement is relative to a named benchmark, information set and evaluation population.

\textbf{Equilibrium and optimization.} An equilibrium comprises feasible plans, optimality, beliefs where needed, budget constraints and all market/resource clearing. The primitives or a stated benchmark define each correspondence \(\mathcal E(p;m)\). Nonemptiness is assumed or proved, never inferred just from compact policy sets. A selected-equilibrium target uses a declared selection rule; a robust target quantifies over all permitted equilibria. Use suprema or infima unless attainment is established. An \(\varepsilon\)-optimal policy is feasible and within \(\varepsilon\) of the finite optimum value. Report an empty feasible set separately.

\textbf{Welfare and accounting.} Normative utility, interpersonal normalization, social weights, numeraire, discounting and the use of public receipts are primitives. Behavioral utility need not coincide with the welfare criterion. Household welfare includes ownership income and rents once; adding profits again to compensating variation generally double-counts them. A monetary equivalent is defined by an explicit transfer and price convention, with monotonicity and range conditions. Average effects, mechanism contributions, transfers and welfare losses are different targets. Infinite sums require convergence and dynamic feasible plans require terminal or no-Ponzi conditions.

\textbf{Source versions.} A working-paper date, revision date and journal publication year are distinguished. A DOI for a journal version is not automatically the DOI of an earlier manuscript. Locators refer to the stated version and printed pages where specified. A metadata-only check does not verify a claim about a full-text conclusion. Every entry's source subsection narrows the provenance claim accordingly.
% END ECONOMICS STATEMENT
\end{document}
